% Fragmento integrable en VII: Relaciones estructurales entre las constantes físicas.
% Insertar después de desarrollos/normalizacion_conjunta.tex.
% Fuentes y cobertura íntegra: editorial/INSERCION_VII.md.
% Sin preámbulo ni macros; usa amsmath y los entornos ya definidos por VII.
\section{Gravity, energy and self-inertia of the joint system}
\label{hmtcuatro:gea:seccion}

The joint normalization allows an operator-level question to be formulated: how geometry, matter and memory enter a single realization. Three links are brought together. The positive character simultaneously preserves its energy, mass and radial readings; variation of the joint transport produces geometric and internal currents of the same energy; and the self-inertial response is expressed through the gravitational radius of that character. Preservation of these links under memory reduction and propagation of the constraints of the coupled action during its regular evolution are proved.

The starting point is the common core of the article. APP produces both sheets with their residues, quotients and incidences; TRIT preserves regime and orientation: $+1$ corresponds to return and memory, $0$ to the present threshold and $-1$ to opening. TPK selects, transports, folds, returns and preserves the enriched state and its history. The prolongation $w_6\to w_{12}\to w_{18}\to w_{24}\to w_{30}\to R_{36}\to G_9$ retains prefix, carry, survival, route and boundary. The dynamics, its nonadic holonomy and the reduced readout follow the order $U_t\to\Gamma_9\to K_{\rm ph}$: phase returns, memory advances and the state is not reset. The five consubstantial constructions of the discrete structure of the continuum are not separated when their readers are realized.

The constants and scales used are outputs of this genealogy. The quadrirelation $(\pi,\varphi,e,\alpha)$ preserves its common origin and its internal constructive order; action, constitutive speed and Barbero--Immirzi are received with their readers. Operator and variational language composes and recognizes these outputs without introducing metrological values to select states or coefficients.

\subsection{Energy, mass and radius of the same character}
\label{hmtcuatro:gea:caracter}

The marked realization of the preceding section is retained:
\[
\Omega=\{(j,k,q,u,v):j,k\in\mathbb Z/12\mathbb Z,\ 
k-j\in\{0,3,6,9\},\ 1\leq q\leq8,\ 
u<v,\ u,v\in\{1,2,4,5,7,8\}\}.
\]
Its cardinality $N_\Omega=12\cdot4\cdot8\binom62=5760$ counts incidences, not matter states or lapses. For the transported inverter $B:V_{\rm in}\to V_{\rm out}$, with $B^*B=BB^*=I/4$, the projector $P=uu^*$ of the common mode has $|u_i|^2=1/N_\Omega$. The archive retains
\[
C_0=(I-P)B,\qquad M_0=PB,\qquad C_0+M_0=B,\qquad\mathsf U=2B.
\]
The images of $C_0,M_0$ are orthogonal and $\mathsf U$ is unitary. The marked expectation $\Delta_\Omega(A)=\sum_iP_iAP_i$, $P_i=e_ie_i^*$, is unique among bimodular expectations onto the diagonal algebra that fix the $P_i$: on a matrix unit, bimodularity requires $\mathcal E(E_{ij})=P_i\mathcal E(E_{ij})P_j$, which is zero if $i\ne j$ and equals $E_{ii}$ if $i=j$.

From $C_0C_0^*=(I-P)/4$ and
$P_iPP_i=P_i/N_\Omega$ follows
\begin{equation}
\Delta_\Omega(C_0C_0^*)=r_\Omega I,\qquad
\Delta_\Omega(M_0M_0^*)=\frac{I}{4N_\Omega},\qquad
r_\Omega=\frac{5759}{23040}.
\label{hmtcuatro:gea:retorno}
\end{equation}
Both readings sum to $I/4$. The longitudinal constitutive rule applies the first coefficient linearly to a length section:
\begin{equation}
L=\alpha^{16}r_\Omega L_*,\qquad D=L\mathsf U.
\label{hmtcuatro:gea:longitud}
\end{equation}
The local norm $\alpha^{16}$ and the section $L_*$ are explicit antecedents; taking an intensity root would be another reader. The radial metric rule is $\|R_Xv\|^2=\|XD^*v\|^2$. For $X=X^*>0$ in the finite fibre, polarization and the unique positive root give $R_X^2=DX^2D^*=L^2\mathsf UX^2\mathsf U^*$ and $R_X=LY$, with $Y=\mathsf UX\mathsf U^*>0$. The spectrum is not selected from a desired gravitational response.

The same section $\hbar>0$ and the conjugate clock $\omega L=c$ determine $E_L=\hbar c/L$ and the readers
\begin{equation}
H_X=E_LY,\qquad M_X=H_X/c^2,\qquad
R_X=LY,\qquad\Lambda_X=LY^{-1}.
\label{hmtcuatro:gea:lectores}
\end{equation}
The notation $H_X$ is reserved for this reader of the character. The joint Hamiltonian $\mathbb H$ introduced later has another domain and is not identified with $H_X$ merely because it shares the name energy.

\begin{proposition}[Common radial coefficient]
\label{hmtcuatro:gea:coeficiente} In the explicit identification of $R_X$ as a reduced gravitational radius, $R_X=GM_X/c^2$, there is a unique compatible coefficient:
\begin{equation}
G=\frac{c^3L^2}{\hbar},\qquad
R_X=\frac{G}{c^4}H_X=\frac{G}{c^2}M_X,\qquad
R_X\Lambda_X=L^2I,\quad H_X\Lambda_X=\hbar cI.
\label{hmtcuatro:gea:radial}
\end{equation}
\end{proposition}
\begin{proof}
Multiplying the readers proves both products. The radial identification is equivalent to $LY=G\hbar Y/(c^3L)$. Cancelling the invertible character determines $G$; back-substitution verifies the equality for every positive character. If two coefficients satisfy it, their difference multiplied by $Y$ is zero and they coincide.
\end{proof}

The longitudinal and metric rules fix the realization; the reduced-radius convention fixes the physical meaning of $G$. The Schwarzschild radius of the same mass is $2R_X$. The equality preserves superpositions and off-diagonal elements in any transported basis, not only mean values. For a differentiable family and a common frame connection $\mathcal A_t$, with $G,c$ fixed,
\begin{equation}
\mathcal D_tR_X=\frac{G}{c^4}\mathcal D_tH_X,\qquad
\mathcal D_tM_X=c^{-2}\mathcal D_tH_X,\qquad
\mathcal D_tO=\dot O+[\mathcal A_t,O].
\label{hmtcuatro:gea:derivada}
\end{equation}
This follows by differentiating the full identity, including the frame. It does not imply that the state is motionless.

For positive modes, or for operators on distinct tensor factors,
\begin{equation}
m(y)=\frac{\hbar}{cL}y,\qquad
\frac{Gm_1m_2}{\hbar c}=y_1y_2,\qquad
\frac{R(y)}{\bar\lambda(y)}=y^2,\quad \bar\lambda(y)=L/y.
\label{hmtcuatro:gea:adimensional}
\end{equation}
Substitution cancels the dimensional sections; operators on distinct factors commute. The dimensionless size depends on the character, without any other free parameter per species. This does not imply a spatial force law or a numerical evaluation without evaluating the character.

The calendar preserves $t_P=L/c$, $t_0=\pi L/(54c)$ and $108t_0=2\pi t_P$. Therefore $L=(54/\pi)\ell_0$, with $\ell_0=ct_0$; the two lengths are not identified. Between sections $\hbar_b=\rho\hbar_a$, $\rho>0$, with $L,c,y_i$ fixed, we have $G_b=G_a/\rho$, $m_{i,b}=\rho m_{i,a}$ and the quantity $Gm_1m_2/(\hbar c)$ remains invariant. Joint sections are compared, not a temporal variation with the other dimensional coordinates artificially held fixed.

\subsection{Static and dynamic memory, and spectral domain}
\label{hmtcuatro:gea:memoria}

In finite dimension, or with bounded coercive blocks, write
\[
Y=\begin{pmatrix}A&F\\F^*&C\end{pmatrix}>0,\quad
Y_{\rm ef}=A-FC^{-1}F^*,\quad \eta=y+C^{-1}F^*b.
\]
Completing the square proves
\begin{equation}
\left\langle\binom b y,Y\binom b y\right\rangle
=\langle b,Y_{\rm ef}b\rangle+\langle\eta,C\eta\rangle.
\label{hmtcuatro:gea:schur-prueba}
\end{equation}
Hence $Y_{\rm ef}>0$, the minimizing extension $y=-C^{-1}F^*b$ and the reconstruction $y=\eta-C^{-1}F^*b$. The residual $\eta$ is retained. Since $\operatorname{Schur}(sY)=sY_{\rm ef}$ for $s>0$, we obtain $H_{X,\rm ef}=E_LY_{\rm ef}$ and $R_{X,\rm ef}=LY_{\rm ef}$. Solving $Y(x,y)=(f,0)$ shows that the boundary block of $Y^{-1}$ is $Y_{\rm ef}^{-1}$, hence $\Lambda_b=LY_{\rm ef}^{-1}$. Direct compression of $Y$ would give $A$: it is not the same operation.

Let $\chi=L/E_L=G/c^4>0$ and $R_X=\chi H_X$. In a decomposition compatible with the domains, we define
\[
\mathcal K_{H_X}(z)=H_{bb}-zI-V(H_{ii}-zI)^{-1}V^*.
\]
\begin{proposition}[Full dynamic memory]
\label{hmtcuatro:gea:resolvente-prop} In the resolvent domain of the interior block,
\begin{equation}
\mathcal K_{R_X}(\chi z)=\chi\mathcal K_{H_X}(z).
\label{hmtcuatro:gea:resolvente}
\end{equation}
If both complements are invertible, their inverses carry the factor $\chi^{-1}$.
\end{proposition}
\begin{proof}
Each radial block is $\chi$ times the energy block and $(\chi H_{ii}-\chi zI)^{-1}=\chi^{-1}(H_{ii}-zI)^{-1}$. Substitution leaves a factor $\chi$ in each term, preserving order. Taking inverses proves the last assertion.
\end{proof}
The energy parameter $z$ is transported to the length $\chi z$; the identity preserves every order of spectral memory. For $|z|<\|H_{ii}^{-1}\|^{-1}$, the resolvent series gives
\[
\mathcal K_{H_X}(z)=H_{\rm est}-zZ+O(z^2),\quad
H_{\rm est}=H_{bb}-VH_{ii}^{-1}V^*,\quad
Z=I+VH_{ii}^{-2}V^*\geq I.
\]
Positivity follows from $Z-I=(H_{ii}^{-1}V^*)^*(H_{ii}^{-1}V^*)$. The static extension $(b,-H_{ii}^{-1}V^*b)$ has squared norm $\langle b,Zb\rangle$: energy and temporal norm come from the same extension.

If $T=Z^{1/2}$ is bounded and invertible, transport is dual:
\begin{equation}
H_{X,c}=T^{-*}H_{X,\rm ef}T^{-1},\quad
R_{X,c}=T^{-*}R_{X,\rm ef}T^{-1},\quad
\Lambda_c=T\Lambda_bT^*.
\label{hmtcuatro:gea:dual}
\end{equation}
Multiplication proves $R_{X,c}=\chi H_{X,c}$, $R_{X,c}\Lambda_c=L^2I$ and $H_{X,c}\Lambda_c=\hbar cI$, without commutativity of $T$ with the character. If $\psi=T(t)b$, $\dot\psi=\dot Tb+T\dot b$ is preserved. The temporal connection term of the generator is not retrospectively incorporated into the radial character without declaring its map.

\begin{proposition}[Spectral extension and compatible refinement]
\label{hmtcuatro:gea:limite} Let $X$ be self-adjoint, positive and injective, and $\mathsf U$ unitary. The readers of \eqref{hmtcuatro:gea:lectores}, defined by spectral calculus of $Y=\mathsf UX\mathsf U^*$, satisfy $R_X=\chi H_X$ on $\operatorname{Dom}(Y)$. Their products with $\Lambda_X$ hold on $\operatorname{Dom}(Y^{-1})$ and have the bounded extensions of \eqref{hmtcuatro:gea:radial}. If isometric embeddings satisfy $X'J=JX$, $\mathsf U'J=K\mathsf U$ and preserve $L,E_L$, then $R'_XK=KR_X$ and $H'_XK=KH_X$ on the transported domains.
\end{proposition}
\begin{proof}
On $P_{\varepsilon,M}=\mathbf1_{[\varepsilon,M]}(Y)$, readers and inverses are bounded and the identities are multiplications of functions of $Y$. For $\psi\in\operatorname{Dom}(Y)$,
\[
\|Y(P_{\varepsilon,M}-I)\psi\|^2
=\int_{(0,\infty)\setminus[\varepsilon,M]}
\lambda^2\,d\langle\psi,E_Y(\lambda)\psi\rangle\longrightarrow0.
\]
Injectivity removes spectral mass at zero and also $P_{\varepsilon,M}\psi\to\psi$. This is convergence in the graph norm. If $\psi\in\operatorname{Dom}(Y^{-1})$, then $Y^{-1}\psi\in\operatorname{Dom}(Y)$ and $YY^{-1}\psi=\psi$, which proves the products and their extensions. Substituting the intertwining relations into $R'_X=L\mathsf U'X'\mathsf U'^*$ proves $R'_XK=KR_X$; the energy calculation is the same.
\end{proof}
The null sector remains separate. This specifies the unbounded extension of the preceding normalization. An arbitrary orthogonal projection is not thereby entitled to the block formulas: unbounded operators require compatible domains or closed forms with controlled blocks.

\subsection{Backreaction in the same transport}
\label{hmtcuatro:gea:retroaccion}

The geometric record $x$, the internal links $u$ and the fields $f$ come from the same enriched state. In an oriented cellularization,
\begin{equation}
E(f,x,u)=\langle r,W(x,u)r\rangle,\quad r=D_Tf,\quad
(D_Tf)_a=f_{t(a)}-T_af_{s(a)},\quad
T_a=U_a^{\rm sp}(x)\otimes R_{\rm int}(u_a).
\label{hmtcuatro:gea:energia-memoria}
\end{equation}
$W=W^*>0$ preserves the blocks between incidences. Spin and internal representation have distinct factors; the chiral realization preserves its blocks and the equivariant Higgs--Yukawa map.

\begin{proposition}[Differential of a single energy]
\label{hmtcuatro:gea:variacion} With fixed fields and $q=Wr$,
\begin{align}
\delta E&=-2\operatorname{Re}\sum_a
\langle q_a,\delta T_af_{s(a)}\rangle+\langle r,\delta Wr\rangle,
\label{hmtcuatro:gea:diferencial}\\
\delta T_a&=\delta U_a^{\rm sp}\otimes R_{\rm int}(u_a)
+U_a^{\rm sp}\otimes\delta R_{\rm int}(u_a).\nonumber
\end{align}
\end{proposition}
\begin{proof}
Differentiating the quadratic form gives two terms conjugate through $W=W^*$ and the term $\langle r,\delta Wr\rangle$. Using $\delta r_a=-\delta T_af_{s(a)}$ proves the first formula; the product rule proves the second. Both variations act on the same $r,q$, with no independently chosen currents.
\end{proof}

\subsubsection{One quantum space and one joint operator}
\label{hmtcuatro:gea:cuantica}

We fix a finite cellularization, a finite set of geometric records $\mathcal X$ and a finite fermionic Fock space $\mathcal F$. The links range over the received compact product $\mathcal G^E$; the doublets $\Phi_v\in\mathbb C^2$ have no amplitude cutoff. The invariant measures $d\mu_x=Z_x^{-1}e^{-S_x}d\mu_0$ have bounded positive densities with bounded inverse and do not depend on $\Phi$ in this chart. The space is
\[
\mathscr H=\bigoplus_{x\in\mathcal X}
L^2(\mathcal G^E\times\mathbb C^{2|V|},
\mu_x\otimes d\Phi;\mathcal F).
\]
The identifications $A_xf=\sqrt{Z_x}e^{S_x/2}f$ are unitary from the common measure and transport the clock to $H_{\rm clk}^{\mu}=A(H_{\rm clk}\otimes I)A^*$. The clock mixes geometric records of the same space. The operator is realized through the form of
\begin{equation}
\mathbb H=H_{\rm clk}^{\mu}+\hbar cK_{\rm gauge}
+d\Gamma(D_T^*WD_T)+H_{\Phi,\rm cin}+V_\Phi
+\widehat H_Y+\widehat H_{\rm tor}.
\label{hmtcuatro:gea:total}
\end{equation}
The Yukawa map is linear in $\Phi,\bar\Phi$ and satisfies $Y(g\Phi)\rho_R(g)=\rho_L(g)Y(\Phi)$. The eliminated contorsion is not simultaneously retained as an independent variable in the base terms of this effective chart.

The scalar-momentum kinetic term has the form $\int\langle\nabla_\Phi\Psi,A(x,u)\nabla_\Phi\Psi\rangle$, with $A$ bounded, uniformly positive and independent of $\Phi$. We require $A(x,u^g)=O_gA(x,u)O_g^{\mathsf T}$ for the real action $O_g$ of the doublets; positivity alone does not prove that covariance. The scalar spatial energy is a positive quadratic potential bounded by $C_{\rm grad}r^2$, $r^2=\sum_v|\Phi_v|^2$, not a bounded operator on all of $L^2$. The volumes satisfy $0<w_{\min}\leq w_v(x)\leq w_{\max}$ and $\lambda_\Phi>0$. The clock, gauge block, transports, $W$ and torsional coefficients are bounded in the chart. The potential retains its absolute value:
\begin{equation}
V_\Phi=\sum_vw_v\lambda_\Phi(|\Phi_v|^2-v_0^2/2)^2
-\sum_vw_v\lambda_\Phi v_0^4/4.
\label{hmtcuatro:gea:potencial}
\end{equation}
The last contribution depends on volume and is not subtracted.

\begin{theorem}[Closed form and joint evolution]
\label{hmtcuatro:gea:forma} The form of \eqref{hmtcuatro:gea:total} is closed and bounded below on $\mathcal Q=\{\Psi\in\mathscr H:\nabla_\Phi\Psi\in L^2,\
r^2\Psi\in L^2\}$. It determines a unique associated self-adjoint operator $\mathbb H$ and the unitary evolution $e^{-it\mathbb H/\hbar}$.
\end{theorem}
\begin{proof}
Cauchy--Schwarz gives $\sum_vw_v\lambda_\Phi|\Phi_v|^4\geq ar^4$, $a=\lambda_\Phi w_{\min}/|V|>0$. Finite Fock space and linear Yukawa give $\|\widehat H_Y(\Phi)\|\leq C_Yr$; torsion is bounded. For $\varepsilon>0$,
\[
C_Yr\leq\varepsilon r^4+
\frac{3C_Y^{4/3}}{4^{4/3}\varepsilon^{1/3}},\qquad
br^2\leq\varepsilon r^4+\frac{b^2}{4\varepsilon}.
\]
The first bound maximizes $C_Yr-\varepsilon r^4$; the second completes the square in $r^2$. Applying them with $\varepsilon=a/4$ to Yukawa and to the negative quadratic term gives
\[
q(\Psi)\geq c_1\|\nabla_\Phi\Psi\|^2+
(a/2)\|r^2\Psi\|^2-C\|\Psi\|^2,\qquad c_1>0.
\]
The norm $\|\Psi\|^2+\|\nabla_\Phi\Psi\|^2+\|r^2\Psi\|^2$ is complete by closedness of the weak derivative and multiplication. These bounds and the upper bounds make the form norm, after adding a constant, equivalent to that complete norm. The representation theorem for closed forms produces the associated operator, and its spectral calculus gives the evolution. Uniqueness refers to this form and domain, not to every alternative formal domain.
\end{proof}

Covariance of $D_T,W$, invariance of measures, the Yukawa identity and internal torsional contraction preserve $q$ and $\mathcal Q$. Uniqueness of the associated operator implies that the compact representation commutes with $\mathbb H$; its Haar average reduces the operator and its evolution. It is not a subsequent corrective average.

The computational cutoff is removed with complete Peter--Weyl blocks and complete isotropic Hermite shells. Regularization and smooth truncation approximate derivatives and the weight $r^2$; bounded densities preserve equivalent norms. Their union is dense in $\mathcal Q$. For the orthogonal projections $P_n$ in the measures $\mu_x$ and the Galerkin operators $\mathbb H_n$,
\[
(\mathbb H_n+\lambda)^{-1}P_nf\longrightarrow
(\mathbb H+\lambda)^{-1}f\qquad(\lambda>C).
\]
The coercive form equation and its Galerkin version give error orthogonality; continuity and coercivity bound it by the best approximation error, which tends to zero. The scalar cutoff and the link-function cutoff of this chart are removed, not the spatial, geometric or fermionic cutoffs. Self-adjointness does not imply that $e^{-t\mathbb H}$ is trace class either.

\subsubsection{Exchange and conservation}
\label{hmtcuatro:gea:intercambio}

For $V=\sum_xP_x\otimes B_x$ and an element $h_{xy}$ of the clock, block multiplication gives
\begin{equation}
[H_{\rm clk}\otimes I,V]_{xy}=h_{xy}(B_y-B_x).
\label{hmtcuatro:gea:conmutador}
\end{equation}
If $h_{xy}\ne0$ and $B_y\ne B_x$, there is geometric--material exchange. If $\mathbb H=A+B+V$, for bounded operators or a common core justifying the derivatives,
\[
\frac d{dt}\langle A\rangle=\frac i\hbar\langle[\mathbb H,A]\rangle,
\quad
\frac d{dt}\langle B\rangle=\frac i\hbar\langle[\mathbb H,B]\rangle,
\quad
\frac d{dt}\langle V\rangle=\frac i\hbar\langle[\mathbb H,V]\rangle.
\]
The sum is zero by $[\mathbb H,\mathbb H]=0$. Energy, including interaction, is conserved, not each sector separately. Geometry belongs to the quantum space; it is not added as an external background. The extended histories retain memory and return; their refinement is not by definition equivalent to spatial refinement.

\subsection{The same action and its currents}
\label{hmtcuatro:gea:accion}

The geometric realization uses a nondegenerate cotetrad $e$, a Lorentz connection $\omega$ and internal connections of the same matter. Let $\Sigma^{IJ}=e^I\wedge e^J$, $P_\gamma=\star+\gamma^{-1}I$, $\gamma\in\mathbb R\setminus\{0\}$ and $\star^2=-I$ on bivectors. The couplings remain fixed during variation. With $\kappa=8\pi G/c^4$, the same action is
\begin{equation}
S_{\rm tot}=\frac{\hbar}{16\pi L^2}
\int\Sigma_{IJ}\wedge(P_\gamma F_\omega)^{IJ}
-\frac{\Lambda\hbar}{8\pi L^2}\int\operatorname{vol}_e+S_m,
\label{hmtcuatro:gea:accion-total}
\end{equation}
where $S_m=S_{\rm YM}+S_\Phi+S_{\rm Weyl}+S_Y$ contracts its indices with the same $e$, which is varied. The coefficients follow from $(2\kappa c)^{-1}=\hbar/(16\pi L^2)$.

With $\delta_\omega S_m=-\tfrac12\int
\delta\omega^{IJ}\wedge\sigma_{IJ}$, $\delta F=D_\omega\delta\omega$ and integration by parts,
\begin{equation}
D_\omega(P_\gamma\Sigma)=\kappa c\sigma.
\label{hmtcuatro:gea:cartan}
\end{equation}
Interior variations or the compatible boundary completion are used. Spin current and metric source are derivatives of the same matter, not two independent inputs.

\subsubsection{Torsional elimination and cross-terms}
\label{hmtcuatro:gea:torsion}

The real inverse is $P_\gamma^{-1}=\gamma^2(-\star+\gamma^{-1}I)/(1+\gamma^2)$: the product of the two polynomials in $\star$ before normalization equals $(1+\gamma^{-2})I$. Denote by $\mathcal Y^{IJ}=\kappa c(P_\gamma^{-1}\sigma)^{IJ}$ the degree-three bivector-valued form, distinct from the character $Y$. For $T^I=D_\omega e^I$, $D_\omega\Sigma^{IJ}=T^I\wedge e^J-e^I\wedge T^J$. With the dual frame $E_I$, its inverse is
\[
\mathcal B^I=\sum_J\iota_{E_J}\mathcal Y^{IJ},\qquad
\vartheta=\tfrac14\sum_I\iota_{E_I}\mathcal B^I,\qquad
T^I=\mathcal B^I-e^I\wedge\vartheta.
\]
Indeed, if $t=\sum_J\iota_{E_J}T^J$, contractions give $\mathcal B^I=T^I+e^I\wedge t$ and then $4t$. The map has zero kernel between spaces of dimension twenty-four, hence is an isomorphism. For $\omega=\mathring\omega(e)+\mathfrak k$,
\[
T_{IJK}=\mathfrak k_{IKJ}-\mathfrak k_{IJK},\qquad
\mathfrak k_{IJK}=\tfrac12(T_{KIJ}-T_{IJK}-T_{JKI}).
\]
The cyclic combination and antisymmetry in the first two indices prove these formulas. For matter affine in $\mathfrak k$, with current independent of it, there is a unique $\mathfrak k_*[e,\sigma]$ linear in the total current $\sigma=\sum_s\sigma_s$.

The curvature expansion preserves the boundary
\[
S_\partial[e,\mathfrak k]=\frac1{2\kappa c}
\int_{\partial M}\Sigma_{IJ}\wedge(P_\gamma\mathfrak k)^{IJ}.
\]
The contorsion-dependent interior density is $\mathcal Q_{e,\gamma}(\mathfrak k)
-\tfrac12\mathfrak k^{IJ}\wedge\sigma_{IJ}$, where $\mathcal Q$ is homogeneous of degree two. Its stationarity evaluated at $\delta\mathfrak k=\mathfrak k_*$ gives $2\mathcal Q(\mathfrak k_*)=
\tfrac12\mathfrak k_*^{IJ}\wedge\sigma_{IJ}$. Therefore,
\begin{equation}
\mathcal L_{\rm spin}^{\rm ef}
=-\tfrac14\mathfrak k_*^{IJ}\wedge\sigma_{IJ}
=-\mathcal Q_{e,\gamma}(\mathfrak k_*).
\label{hmtcuatro:gea:spin}
\end{equation}
Linearity preserves cross-terms between species. Contorsion is eliminated once; its eliminated linear interaction and the effective quartic term are not added simultaneously.

In the spinorial specialization, the received coefficient $C_\gamma=3\kappa c\hbar^2\gamma^2/[16(1+\gamma^2)]$ gives
\begin{equation}
\frac{C_\gamma}{\hbar}
=\frac{3\pi}{2}L^2\frac{\gamma^2}{1+\gamma^2}.
\label{hmtcuatro:gea:coeficiente-torsion}
\end{equation}
It suffices to substitute $\kappa=8\pi G/c^4$ and $\hbar G=c^3L^2$, without changing the branch of $\gamma$. In the Fock space of the chart, the total current preserves the contraction
\[
\widehat Q_{{\rm tot},v}=\sum_{a=1}^4s_a
\{d\Gamma(\widetilde M_{a,v})^2-d\Gamma(\widetilde M_{a,v}^2)\},
\quad s=(-1,-1,-1,+1),
\]
\[
\widehat H_{\rm tor}
=-\sum_v\frac{N_vcC_\gamma}{\mathcal V_v}\widehat Q_{{\rm tot},v}.
\]
The matrices $\widetilde M_{a,v}$ represent the current components. $N_v$ is the lapse, distinct from $N_\Omega$ and from an occupation. The normal-ordering identity removes the one-particle contraction, not the cross-terms between species. The signature remains in $s_a$.

\subsubsection{Metric source and Legendre transform of the same action}
\label{hmtcuatro:gea:legendre}

Stationary elimination preserves variation of the other variables: the chain rule contains a term in $\delta\mathfrak k_*$ multiplied by the connection equation, which is zero. The reduced action is
\[
S_{\rm red}=\frac1{2\kappa c}\int(R[g]-2\Lambda)\operatorname{vol}_e
+S_m^{\rm LC}+S_{\rm spin}^{\rm ef}+S_\partial.
\]
We define the action source by $\delta_g(S_m^{\rm LC}+S_{\rm spin}^{\rm ef})
=-\tfrac12\int\mathcal T^S_{\mu\nu}\delta g^{\mu\nu}
\operatorname{vol}_e$. The interior gravitational variation produces
\begin{equation}
G_{\mu\nu}+\Lambda g_{\mu\nu}
=\kappa c\mathcal T^S_{\mu\nu}
=\kappa(T^{\rm LC}_{\mu\nu}+U^{\rm spin}_{\mu\nu}),
\quad T^{\rm LC}+U^{\rm spin}=c\mathcal T^S.
\label{hmtcuatro:gea:metrica}
\end{equation}
The volume variation of the constant term of the potential is included. The factor $c$ corresponds to the cotetrad on the length line; action and physical sources are not mixed.

With $a=(2\kappa c)^{-1}$ and $\mathcal B_{IJ}=(P_\gamma\Sigma)_{IJ}$, the geometric potential is $\theta_{\rm geom}=a\mathcal B_{IJ}\wedge\delta\omega^{IJ}$. Decomposing $\omega=\omega_0dt+\underline\omega$ and $F_{0a}=\dot\omega_a-D_a\omega_0$ identifies the kinetic term; integration by parts preserves Gauss and the boundary. With $e_0=Nn+N^ae_a$ and the Legendre transform of the same matter,
\begin{equation}
S_{\rm tot}=\int dt\,[\Theta_{\rm red}(\dot z)
-\mathcal C[N]-\mathcal V[N^a]-\mathcal G_L[\omega_0]
-\mathcal G_{\rm int}[A_0]]+S_\partial.
\label{hmtcuatro:gea:canonica}
\end{equation}
Lapse, shift and gauge retain their degenerate directions. Second-class constraints are reduced before using $\Theta_{\rm red}$ or are expressly retained. First-order spinors are not assigned a fictitious velocity inverse.

For contorsion and its momentum, the constraints are $\pi_{\mathfrak k}=0$ and $\partial\mathcal Q/\partial\mathfrak k-\sigma/2=0$. Their matrix has blocks $\left(\begin{smallmatrix}0&-M\\M^{\mathsf T}&B\end{smallmatrix}\right)$, with $M$ invertible by the preceding reconstruction; it is therefore invertible. The Dirac bracket performs this elimination and separately preserves the graded fermionic constraints. In a regular bosonic direction,
\[
H_\Phi=N\left(\frac{p^2}{2w}+wV+H_{\rm grad}\right),\quad
p=\frac wN D_t\phi,\quad
L_\Phi=\frac{w}{2N}|D_t\phi|^2-NwV-NH_{\rm grad}.
\]
The volume $w=\sqrt{\det q}$ is not frozen. For $K_{ij}=(\dot q_{ij}-\mathcal L_{N^a}q_{ij})/(2N)$,
\[
\Pi^{ij}=\frac{\sqrt q}{2\kappa c}(K^{ij}-q^{ij}K),\qquad
K_{ij}=\frac{2\kappa c}{\sqrt q}
(\Pi_{ij}-\tfrac12q_{ij}\Pi).
\]
The trace and the traceless part prove the inverse. If the total momentum contains a material frame contribution, $\Pi=p-p_m$ is used, preserving $p_m$. The same action and its source are thereby recovered. The Legendre transform does not by itself prove that an arbitrary quantization of \eqref{hmtcuatro:gea:canonica} coincides with \eqref{hmtcuatro:gea:total}.

\subsection{Constraint propagation with the total source}
\label{hmtcuatro:gea:restricciones}

Let $E_{\mu\nu}=G_{\mu\nu}+\Lambda g_{\mu\nu}
-\kappa c\mathcal T^S_{\mu\nu}$. On the matter equations, diffeomorphism invariance of the reduced functional implies $\nabla^\mu\mathcal T^S_{\mu\nu}=0$: an interior variation generated by $\xi$ leaves the integral of $\mathcal T^S_{\mu\nu}\nabla^\mu\xi^\nu$, while the matter equations cancel the other terms. Integrating by parts and varying $\xi$ proves zero divergence, before imposing $E=0$. Bianchi and the fixed couplings give $\nabla_\mu E^{\mu\nu}=0$.

\begin{proposition}[Preservation of constraint data]
\label{hmtcuatro:gea:propagacion} A regular solution of the matter and metric evolution equations preserves the initially satisfied constraints throughout its interval of regular existence.
\end{proposition}
\begin{proof}
In Gaussian coordinates $g=-d\tau^2+q_{ij}(\tau,x)dx^idx^j$, we impose $E_{ij}=0$ and write $C=E_{00}$, $C_i=E_{0i}$, $K_{ij}=\dot q_{ij}/2$ and $\mathcal K=q^{ij}K_{ij}$. We have $E^{00}=C$, $E^{0i}=-C^i$, $E^{ij}=0$, $\Gamma^i{}_{0j}=K^i{}_j$ and $\Gamma^\mu{}_{\mu0}=\mathcal K$. The temporal and spatial components of the divergence give
\[
\partial_\tau C-D_iC^i+\mathcal KC=0,\qquad
\partial_\tau C^i+\mathcal KC^i+2K^i{}_jC^j=0.
\]
On lowering the index, $\dot q_{ij}=2K_{ij}$ cancels the last spatial term. Exactly the following remain
\begin{equation}
\partial_\tau C-D_iC^i+\mathcal KC=0,\qquad
\partial_\tau C_i+\mathcal KC_i=0.
\label{hmtcuatro:gea:propagacion-sistema}
\end{equation}
The second homogeneous equation maintains $C_i=0$ and the first becomes $\dot C+\mathcal KC=0$, preserving $C=0$. $q$ and $\mathcal K$ are those of the coupled solution, not a prescribed background.
\end{proof}
The Gaussian chart locally proves a tensorial assertion, without freezing $q^{-1}$ or replacing the source. Global singularity-free existence is not asserted for every datum. Nor is variational propagation by itself a quantum commutator identity for arbitrary normal deformations.

\subsection{Boundary and self-inertial operator}
\label{hmtcuatro:gea:autoinercia}

Let $b:\mathcal X_\infty^{\rm enr}\to B_{\rm ef}$, $B_{\rm ef}=\operatorname{im}(b)$. A response $I_v$ factors uniquely as $I_v=\overline I_v\circ b$ if and only if it is constant on the fibres of $b$. Necessity is immediate; for sufficiency, $\overline I_v(\beta)=I_v(x)$ with $b(x)=\beta$ is independent of the representative. The effective image gives uniqueness. Strong Machian dependence adds states with the same local restriction and different responses.

The hinge of the same state produces $r_{\rm av}=\pi/\varphi^2$ and $r_{\rm av}^J=\varphi/\pi^2$. In the common-amplitude sector, $\mathcal L_{\rm ar}=\pi\mu_{\rm ar}\mathcal M(x)$ and $\mathcal L_{\rm vol}=\mu_{\rm vol}\mathcal M(x)$, $\mathcal M(x)>0$, with $\mu_{\rm ar}/\mu_{\rm vol}=\varphi^{-2}$. Dividing by $\mathcal L_{\rm vol}$ gives readers of the same type $A_\partial=r_{\rm av}$ and $V_\partial=1$. The quotient cancels the amplitude, not the history preserved in the archive.

For an orientation $a$, $q_a(x)=h_a(x)/(108t_0(x))=\hbar_a(x)/t_P(x)$ descends to $Q:B_{\rm ef}\to\mathcal L_E^+$ if and only if $b(x)=b(x')$ implies $h_a(x)/t_0(x)=h_a(x')/t_0(x')$. This is the factorization criterion applied to the clock. It is satisfied by fixed charts and by variable charts whose quotient descends. In the radial chart, $Q=\hbar_ac/L$. With energies on that same line, $E_1E_2/Q^2$ is dimensionless and basis-independent.

For complementary orthogonal projections,
\begin{align}
w_A&=\frac{A_\partial r_{\rm av}}
{A_\partial r_{\rm av}+V_\partial r_{\rm av}^J}
=\frac{\pi^4}{\pi^4+\varphi^5},\qquad w_V=1-w_A,
\label{hmtcuatro:gea:pesos}\\
W_{\rm av}&=w_AP_A+w_VP_V,\qquad
\overline I_v=\frac{E_1E_2}{Q^2}W_{\rm av}.
\label{hmtcuatro:gea:respuesta}
\end{align}
Multiplying the numerator and denominator of $r_{\rm av}^2/(r_{\rm av}^2+r_{\rm av}^J)$ by $\varphi^4\pi^2$ proves the evaluation. The distinction between areal reader and weighting is preserved: the regional mark is generated once. $0<w_A,w_V<1$ makes $W_{\rm av}$ positive and invertible, and
\[
\langle z,\overline I_vz\rangle
=\frac{E_1E_2}{Q^2}
(w_A\|P_Az\|^2+w_V\|P_Vz\|^2)>0\quad(z\ne0).
\]
Fully exchanging $(A_\partial,r_{\rm av},P_A)$ and $(V_\partial,r_{\rm av}^J,P_V)$ permutes the summands. Unitary transport conjugates the response and preserves its spectrum; implementing the exchange within one space requires equal sector dimensions. A common rescaling of $A_\partial,V_\partial$ does not change the weights and does not prove global dependence. With fixed energies and weights, $Q\mapsto qQ$ changes the response by $q^{-2}$; the strong witness also preserves the same local restriction.

\begin{theorem}[Self-inertia and radius of the same character]
\label{hmtcuatro:gea:autoinercial-radial} With the same clock and orientation,
\begin{equation}
\overline I_v=\frac{G_aE_1E_2}{\hbar_ac^5}W_{\rm av}
=\frac{G_am_1m_2}{\hbar_ac}W_{\rm av}.
\label{hmtcuatro:gea:mach-gravedad}
\end{equation}
For a probe $y_p>0$, $E_p=Qy_p$ and $\bar\lambda_p=L/y_p$, the spectral extension satisfies
\begin{equation}
\widehat I_{p,Y}=y_pY\otimes W_{\rm av}
=\frac{E_p}{Q^2}H_X\otimes W_{\rm av}
=\frac{R_X}{\bar\lambda_p}\otimes W_{\rm av}
=\frac{G_am_p}{\hbar_ac}M_X\otimes W_{\rm av}.
\label{hmtcuatro:gea:identidad-autoinercial}
\end{equation}
\end{theorem}
\begin{proof}
$Q^2=\hbar_a^2c^2/L^2$ and $\hbar_aG_a=c^3L^2$ give $Q^{-2}=G_a/(\hbar_ac^5)$; $E_i=m_ic^2$ is then used. For $Y=\sum_jy_jP_j$, applying the reader to each energy $Qy_j$ produces $\sum_jy_py_jP_j\otimes W_{\rm av}$, independently of the eigenbasis. Spectral calculus extends the linear reader to the general positive self-adjoint character. Decomposition by $P_A,P_V$ gives the operators $y_pw_AY$ and $y_pw_VY$, with their multiplicities: they are positive self-adjoint operators on the domain of $Y\otimes I$. Substituting the radial readers proves the other expressions.
\end{proof}

If an additional realization effectively identifies its full operator with $H_X=QY$, it preserves that identification and its domains. For $H_X=\sum_rH_r$ on a common core, $\widehat I_{p,Y}=(E_p/Q^2)\sum_rH_r\otimes W_{\rm av}$. Positivity of each summand is not required; positivity of the character of the identified operator is. No terms are suppressed or energy origin changed to impose it. The theorem does not automatically assign this reader to $\mathbb H$.

In the blocks of \eqref{hmtcuatro:gea:schur-prueba}, the interior inverse of $y_pY\otimes W_{\rm av}$ is $y_p^{-1}C^{-1}\otimes W_{\rm av}^{-1}$. Multiplication gives
\begin{equation}
\operatorname{Schur}(\widehat I_{p,Y})
=y_pY_{\rm ef}\otimes W_{\rm av}
=\frac{R_{X,\rm ef}}{\bar\lambda_p}\otimes W_{\rm av}.
\label{hmtcuatro:gea:schur-autoinercia}
\end{equation}
The cross-blocks are not deleted. For $a_s=y_pw_s>0$, $s=A,V$, the same calculation gives $\mathcal K_{a_sY}(a_sz)=a_s\mathcal K_Y(z)$ in the interior resolvent domain, preserving full memory. If $Y'K=KY$ and $W'_{\rm av}J=JW_{\rm av}$, with $y_p$ and the intensive section fixed,
\begin{equation}
\widehat I'_{p,Y'}(K\otimes J)
=(K\otimes J)\widehat I_{p,Y}.
\label{hmtcuatro:gea:naturalidad}
\end{equation}
This is proved by substitution. The spectral cutoffs of Proposition \ref{hmtcuatro:gea:limite} converge in graph norm; the two fixed weights transmit convergence to the response. Naturality is that of these effective intertwiners, not that of a different family of spatial maps identified by name alone.

\subsection{Self-inertia of the state and acceleration of its projection}
\label{hmtcuatro:gea:proyeccion}

In a differentiable realization, the joint state $\Gamma=(x,\mathcal O,m)$ has connection $\nabla^{\rm tot}$ on $\mathcal N$. A smooth projection $\pi_S:\mathcal N\to\mathcal N_S$, with connection $\nabla^S$, preserves the transports and readers it realizes. The observed frame is $\mathfrak F_{\mathcal O}(\Gamma)
=\operatorname{Res}_{\mathcal O}(\Gamma,\nabla^{\rm tot})$. For a twice differentiable curve and $\gamma_S=\pi_S\circ\Gamma$, self-inertia means $\nabla^{\rm tot}_{\dot\Gamma}\dot\Gamma=0$. The defect is
\[
\mathcal K_{\pi_S}(\dot\Gamma,\dot\Gamma)
=\nabla^S_{\dot\gamma_S}\dot\gamma_S
-d\pi_S(\nabla^{\rm tot}_{\dot\Gamma}\dot\Gamma).
\]
\begin{proposition}[Readout acceleration]
\label{hmtcuatro:gea:aceleracion}
For a self-inertial trajectory,
\begin{equation}
a_S=\mathcal K_{\pi_S}(\dot\Gamma,\dot\Gamma),\qquad
(\mathcal K_{\pi_S})^a{}_{BC}
=\partial_B\partial_C\pi_S^a
+(\Gamma^S)^a{}_{bc}\partial_B\pi_S^b\partial_C\pi_S^c
-\partial_D\pi_S^a(\Gamma^{\rm tot})^D{}_{BC}.
\label{hmtcuatro:gea:defecto}
\end{equation}
\end{proposition}
\begin{proof}
The chain rule gives $\ddot\gamma_S^a=\partial_B\pi_S^a\ddot\Gamma^B
+\partial_B\partial_C\pi_S^a\dot\Gamma^B\dot\Gamma^C$. Adding the subsystem connection and subtracting the image of the total acceleration cancels the terms in $\ddot\Gamma$. The remaining terms are the displayed tensor contracted with the velocities. Self-inertia sets the total acceleration to zero and leaves the equality.
\end{proof}
The geometric response factors through $b$ when the connection, projection and kinematic datum descend to the boundary, or the latter is fixed in the comparison. If velocity changes, it is retained as an argument. The self-inertial whole can thus yield accelerated subsystems as readouts without an external absolute frame.

Self-inertia does not by itself cancel curvature or holonomy. A datum transported around a loop returns if and only if it lies in the fixed subspace of its holonomy. Return of the full state additionally requires an injective realization in the sector being compared: visible return and return of history remain distinct.

\subsection{Result and scope checks}
\label{hmtcuatro:gea:conclusion}

The same character fixes mass, energy and radius; their proportionality preserves static and spectral memory and dual normalization. Geometric and internal transports enter the differential of a single energy. The joint action produces the geometric and torsional sources and preserves the constraints on its regular solution. The self-inertial operator is the radius of the character divided by the conjugate length of the probe, while acceleration is the projection defect of the state.

The last two objects have distinct types: $\widehat I_{p,Y}$ is identified neither with $\mathcal K_{\pi_S}$ nor with the Einstein tensor. Nor is the proved incorporation replaced by the different requirement of cancelling a total quantum curvature for arbitrary normal deformations. That statement has its own generators and domains; it is neither asserted here nor declared absent from the corpus because it is not this conclusion.

The falsifiers are concrete. Changing only $G$ or $\hbar$ with $L,Y$ fixed breaks \eqref{hmtcuatro:gea:radial}; replacing $L$ with $\ell_0$ without $54/\pi$ alters the clock; deleting $F$ modifies the reduction; jointly rescaling the areal and volumetric readers does not prove Machian separation; a map that does not intertwine $Y$ does not satisfy \eqref{hmtcuatro:gea:naturalidad}. The rational tests are reproducible finite checks, not target values of the generator or substitutes for the domain and limit proofs.
